\paragraph{Grid-following converter.} A converter exports 1 pu to an infinite bus through $X_g=1/\mathrm{SCR}$ in a quasi-static phasor model, with the SCR swept from 10 down to 1.0. The GFL converter uses an SRF-PLL (damping 0.707, bandwidth 10, 20 or 40 Hz), an AC-voltage PI loop holding $|v_{pcc}|=1$ and a 2 ms current-loop lag; the GFM converter is a VSM with $X_v=0.15$ pu, $D=20$ pu and $H=2$ or 5 s. Fig.~\ref{fig:C1_weakgrid_stability}(a) shows that the GFL converter loses small-signal stability before it runs out of equilibria: an oscillatory eigenvalue pair crosses into the right half plane at a critical SCR of 1.367, 1.332 and 1.298 for the 10, 20 and 40 Hz PLL, against static limits of SCR 2.0 (unity power factor, $i_q=0$) and 1.0 (with voltage support). Inside the IEEE Std 1204 weak-grid band (SCR below 3) the least-damped oscillatory mode of the 10 Hz PLL falls from a damping ratio of 0.702 at SCR 10 to 0.195 at SCR 1.5. Below SCR 2 a coupled PLL and voltage-loop oscillation of about 5 to 7 Hz takes over and becomes unstable.

\paragraph{What this model does not show.} In this quasi-static model the PLL bandwidth shifts the GFL critical SCR only slightly (1.37, 1.33 and 1.30, a faster PLL being marginally better). The well-known penalty of a fast PLL needs the electromagnetic current-loop and network dynamics that this model omits; these are planned for Step 2, and we do not claim the penalty here.

\paragraph{Grid-forming converter and the duality.} Fig.~\ref{fig:C1_weakgrid_stability}(b) and the eigenvalue loci of Fig.~\ref{fig:C2_eig_loci} show that the VSM stays stable at every SCR down to its static limit (last equilibrium at SCR 1.01) for both $H=2$ s and $H=5$ s, and that its swing-mode damping rises as the grid weakens: from 0.090 at SCR 10 to 0.194 at SCR 1.5 for $H=5$ s (0.142 to 0.307 for $H=2$ s). This is the duality of Li, Gu and Green \cite{li}: grid-following is robust in strong grids and fragile in weak ones, grid-forming is the opposite, and its weak spot is a lightly damped swing mode in strong grids ($\zeta\approx0.09$ at SCR 10 for $H=5$ s), which a larger damping gain or a virtual-impedance damper addresses. The result matters for India because large renewable complexes, such as those in western Rajasthan, feed long EHV corridors far from synchronous generation, which is the low-SCR regime. No substation-specific SCR is asserted here; site values must come from CEA and CTUIL system studies. The caveats are those of the model: a quasi-static network without line dynamics, no current-loop PI detail, no converter current limit, and a constant internal voltage of the VSM in the small-signal analysis.
